\documentclass[a4paper,11pt]{amsart}
\usepackage[margin=1in]{geometry}
\usepackage[T1]{fontenc}
\usepackage{lmodern}
\usepackage{microtype}
\usepackage{amsmath,amssymb,mathtools,amsthm}
\usepackage{booktabs,tabularx,array}
\usepackage{enumitem}
\usepackage[colorlinks=true,linkcolor=blue,citecolor=blue,urlcolor=blue]{hyperref}
\usepackage[capitalize,nameinlink,noabbrev]{cleveref}

\allowdisplaybreaks[2]
\emergencystretch=3em
\setlist[enumerate]{leftmargin=*,itemsep=2pt,topsep=4pt}
\setlist[itemize]{leftmargin=*,itemsep=2pt,topsep=4pt}

\newtheorem{theorem}{Theorem}[section]
\newtheorem{proposition}[theorem]{Proposition}
\newtheorem{lemma}[theorem]{Lemma}
\newtheorem{corollary}[theorem]{Corollary}
\newtheorem{conjecture}[theorem]{Conjecture}
\newtheorem{metaconjecture}[theorem]{Meta-conjecture}
\newtheorem{problem}[theorem]{Problem}
\newtheorem{definition}[theorem]{Definition}
\theoremstyle{remark}
\newtheorem{remark}[theorem]{Remark}

\crefname{theorem}{theorem}{theorems}
\Crefname{theorem}{Theorem}{Theorems}
\crefname{proposition}{proposition}{propositions}
\Crefname{proposition}{Proposition}{Propositions}
\crefname{lemma}{lemma}{lemmas}
\Crefname{lemma}{Lemma}{Lemmas}
\crefname{corollary}{corollary}{corollaries}
\Crefname{corollary}{Corollary}{Corollaries}
\crefname{conjecture}{conjecture}{conjectures}
\Crefname{conjecture}{Conjecture}{Conjectures}
\crefname{metaconjecture}{meta-conjecture}{meta-conjectures}
\Crefname{metaconjecture}{Meta-conjecture}{Meta-conjectures}
\crefname{problem}{problem}{problems}
\Crefname{problem}{Problem}{Problems}
\crefname{definition}{definition}{definitions}
\Crefname{definition}{Definition}{Definitions}
\crefname{remark}{remark}{remarks}
\Crefname{remark}{Remark}{Remarks}

\DeclareMathOperator{\RS}{RS}
\DeclareMathOperator{\MCA}{MCA}
\DeclareMathOperator{\im}{im}
\newcommand{\F}{\mathbb F}
\newcommand{\B}{\mathbb B}
\newcommand{\Z}{\mathbb Z}
\newcommand{\Gam}{\Gamma}
\newcommand{\Om}{\Omega}
\newcommand{\Lam}{\Lambda}
\newcommand{\calE}{\mathfrak E}
\newcommand{\Htwo}{H_2}
\newcommand{\barN}{\overline N}
\newcommand{\abs}[1]{\lvert#1\rvert}
\newcommand{\floor}[1]{\left\lfloor#1\right\rfloor}
\newcommand{\ceil}[1]{\left\lceil#1\right\rceil}
\newcommand{\defeq}{\mathrel{:=}}

\title{Conjectures and Barriers for RS--MCA}
\author{Przemek Chojecki}
\address{Ulam.ai}
\email{prz.chojecki@gmail.com}
\urladdr{https://github.com/przchojecki/rs-mca}
\date{July 19, 2026}
\hypersetup{
  pdftitle={Conjectures and Barriers for RS--MCA},
  pdfauthor={Przemek Chojecki},
  pdfsubject={Profile conjectures, inverse problems, and finite barriers for Reed--Solomon mutual correlated agreement},
  pdfkeywords={Reed-Solomon codes, mutual correlated agreement, Proximity Prize, profile compiler, additive energy, Balog-Szemeredi-Gowers theorem}}
\subjclass[2020]{94B35, 11T71, 05B40, 11B30}
\keywords{Reed--Solomon codes, mutual correlated agreement, profile certificates,
additive energy, smooth domains, proximity gaps}

\begin{document}
\raggedbottom

\begin{abstract}
As a companion to the proved bounds in \emph{Shortening Bounds for Reed--Solomon MCA}, we formulate the complete finite Reed--Solomon mutual correlated agreement (MCA) problem and a pole-aware conjectural positive-density profile envelope.  An exact first-match compiler and realized-image moment and incidence inequalities isolate the payments required for a safe certificate without confusing supports, pairs, rays, and affine slopes.  We prove a ceiling-normalized moment obstruction, a projective incidence theorem and benchmark dimension diagnostic, and a moving-scale consequence of the previously proved implication from a Sidon payment through the Balog--Szemer\'edi--Gowers theorem and Boolean-cube growth.  Exact unsafe edges, repository audits, and numerical margins lead to four direct adjacent conjectures---two MCA and two auxiliary list inequalities---whose safe sides still require primitive-fiber, residual-projection, algebraic-routing, and add-back payments.  In a collision-nonbinding, subexponential-budget identity-candidate branch, we conjecture a non-oracular exhaustive atlas whose exact unsafe--safe bracket, through a proved crossing reduction, yields \(\delta^*_{C_n,\mathrm{off,sup}}=1-\rho_n-g^*(\rho_n,\log_2\abs{\mathbb B_n})+o(g_n^*)\); the general identity lower route instead uses an exact pole-adjusted target, and no matching bracket is claimed for the unrestricted smooth or circle problem.
\end{abstract}

\maketitle

\section{Introduction}\label{sec:intro}

The preliminary Proximity Prize asks one to determine the complete security
frontier of
Reed--Solomon mutual correlated agreement (MCA) on structured evaluation
domains \cite{ProximityPrize,ABF26}.  Its grand MCA challenge takes
\[
 C=\RS_{\F}(D,k)
 =\{(f(x))_{x\in D}:f\in\F[X],\ \deg f<k\},\qquad n=\abs D,
\]
at a rate \(k/n\in\{1/2,1/4,1/8,1/16\}\).  The preliminary statement phrases
the task as finding the largest safe radius in the official interval
\([0,1]\) at a target such as \(\varepsilon^*=2^{-128}\).  That wording hides a finite staircase,
and the soundness gap between an attack and a safe certificate is meaningful
only after its integer and endpoint conventions have been fixed.

For a received pair \(r=(r_0,r_1)\), the affine received line is
\(\{r_0+\gamma r_1:\gamma\in\F\}\).  At agreement \(a\), a slope
\(\gamma\) is MCA-bad if there are a support \(S\subseteq D\), with
\(\abs S\ge a\), and \(h\in\F[X]_{<k}\) for which
\((r_0+\gamma r_1)|_S=h|_S\), while \(r_0\) and \(r_1\) are not both
explained by degree-less-than-\(k\) polynomials on that same support.  Thus
MCA counts distinct affine slopes, not supports, locators, polynomials,
pairs, or projective directions.  For a nonempty challenge set
\(\Gam\subseteq\F\), let \(B_{C,\Gam}^{\MCA}(a)\) be the maximum number of
bad slopes in \(\Gam\), over all received pairs.  At radius \(\delta\), put
\[
 a(\delta)=n-\floor{\delta n},\qquad
 \varepsilon_{\MCA,\Gam}(C,\delta)
 =\frac{B_{C,\Gam}^{\MCA}(a(\delta))}{\abs\Gam}.     \tag{1.1}\label{eq:intro-error}
\]
The official challenge is the full-field case \(\Gam=\F\).  The
challenge-restricted notation is our refinement for sampled variants; it
does not change the support-wise witness predicate.  We abbreviate
\(B_C^{\MCA}=B_{C,\F}^{\MCA}\).

\begin{problem}[Complete finite Proximity Prize problem]\label{prob:finite-prize}
Given the official data above, determine the complete safe set
\[
 \mathsf S_{C,\Gam}(\varepsilon^*)
 =\{\delta\in[0,1]:
       B_{C,\Gam}^{\MCA}(a(\delta))
       \le B^*\},\qquad
 B^*=\floor{\varepsilon^*\abs\Gam},                \tag{1.2}\label{eq:intro-safe-set}
\]
and not only a decimal approximation to one boundary point.  When the safe
set is nonempty, report its real supremum
\(\delta_{C,\Gam,\mathrm{off,sup}}^*\), its largest safe grid point
\(\delta_{C,\Gam,\mathrm{off,grid}}^*\in\{0,1/n,\ldots,1\}\), and whether
the supremum is attained.  When it is empty, report that fact separately;
the number zero may be used as a placeholder but must not be read as saying
that radius zero is safe.
The concrete box emphasized in the accompanying preliminary problem
discussion includes \(k\le2^{40}\), \(\abs\F<2^{256}\), and target
\(2^{-128}\) \cite{ABF26}; the definitions and endpoint bookkeeping here do
not depend on those numerical caps.
\end{problem}

The complete answer has three branches.  For a proper MDS code,
\cref{prop:staircase} proves that \(B_{C,\Gam}^{\MCA}(n)=1\) and that the
numerator is constant for \(a\le k+1\).  Hence \(B^*=0\) gives an empty safe
set.  If \(B_{C,\Gam}^{\MCA}(k+1)\le B^*\), the row is saturated:
\(\mathsf S_{C,\Gam}=[0,1]\), both official thresholds equal \(1\), and the
endpoint is safe.  Otherwise \(B^*\ge1\), and if
\[
 a_C^*=\min\{a\in\{k+1,\ldots,n\}:
                   B_{C,\Gam}^{\MCA}(a)\le B^*\},   \tag{1.3}\label{eq:intro-first-safe}
\]
then
\[
 \mathsf S_{C,\Gam}
 =\left[0,\frac{n-a_C^*+1}{n}\right),\quad
 \delta_{C,\Gam,\mathrm{off,grid}}^*=\frac{n-a_C^*}{n},\quad
 \delta_{C,\Gam,\mathrm{off,sup}}^*=\frac{n-a_C^*+1}{n}, \tag{1.4}\label{eq:intro-staircase}
\]
and the supremal endpoint is unsafe.  These empty, saturated, grid, and
half-open cases are part of the finite Prize answer; a subcapacity threshold
alone is incomplete.

The Prize statement uses a smooth evaluation domain.  Our refined term
\emph{admissibly smooth} means that the row data include specified complete
equal-fiber polynomial folding maps.  Here a folding \(\pi:D\to D_\pi\) is
complete and \(c\)-to-one, for an integer \(c\ge1\), when
\(\pi(D)=D_\pi\) and every fiber
\(\pi^{-1}(y)\), \(y\in D_\pi\), has exactly \(c\) elements.  The principal examples are
positive-density multiplicative domains, with power maps \(x\mapsto x^c\),
and Chebyshev twin-coset domains, with the corresponding maps \(T_c\).  We
write \(D\subseteq\B\subseteq\F\) when this locator and folding algebra is
defined over a designated coefficient-and-folding field \(\B\), while slopes
and codewords live over \(\F\).  This field is part of the public row data;
for the asymptotic prediction it is the field generated by \(D\) and the
coefficients of the declared folding maps, not an arbitrarily enlarged
intermediate field.  A circle presentation enters only through an explicit
coordinate bijection and nonzero diagonal scaling that give an
agreement-preserving monomial equivalence with such an RS row
\cite{HLP24,ChoShortening}.  Geometric resemblance alone is not used.
For a sequence, \emph{positive density} means
\(\liminf_n\abs{D_n}/\abs{\B_n}>0\); hence
\(\log_2\abs{\B_n}=\log_2n+O(1)=o(n)\).

The proved literature does not yet determine the grand frontier.  BCIKS
established the foundational Reed--Solomon proximity-gap framework
\cite{BCIKS20}.  Hab\"ock gave the quadratic predecessor of a Johnson-range
MCA estimate \cite{Habock25}; BCHKS state a linear-in-\(n\) Johnson-range
estimate, which the companion theorem paper treats as an imported input,
as well as complementary proximity-gap obstructions \cite{BCHKS25}.  Jo
proves an all-test-size MDS circuit-incidence envelope and an exact
agreement-set shortening theorem, yielding polynomial control at any fixed
number of steps beyond Johnson at fixed reciprocal-integer rates
\cite{Jo26}.  BCHKS use degree at most \(d\); for \(\RS_{<K}\), their strict
Johnson hypothesis is \((N-E)^2>N(K-1)\).

Recent attacks show why no identity-only capacity law can be universal.
Crites and Stewart disprove several correlated-agreement and MCA
up-to-capacity conjectures \cite{CritesStewart25}.  Krachun, Kazanin, and
Hab\"ock construct multiplicative-subgroup rows over prime fields with
\(2^{\Omega_\rho(1/\eta)}\) close points on one affine line at a gap
\(\eta=\Theta_\rho(1/\log n)\) from capacity \cite{KKH26}.  Gao, Yang, Xu,
and Kan give a constructive conversion from list-decoding counterexamples to
MCA lower bounds for a related code, including an RS specialization
\cite{GaoYangXuKan26}.  These are nonidentity lower mechanisms: they sharpen
the attack side and the profile taxonomy, but do not furnish fixed-row safe
certificates for the benchmarks below.

The preceding theorem paper, \emph{Shortening Bounds for Reed--Solomon MCA},
version~9.2 \cite{ChoShortening}, is the source of the unconditional
RS--MCA results used here.  It proves the official MDS plateau and exact
endpoint conversion, the endpoint-exact linear-budget staircase in its
below-half-distance range, finite MDS circuit and shortening certificates
beyond Johnson, and a one-sided positive-radius safe-frontier certificate
for positive exponential budgets.  It also proves the exact locator-prefix
lower construction with collision-aware pole conversion, structured-domain
and circle transfers under proved monomial equivalence, and the implication
from an image-normalized Sidon payment to primitive Boolean-fiber flatness.
The latter implication uses Balog--Szemer\'edi--Gowers and Boolean-cube
difference growth; its Sidon premise and residual ray compiler remain open.
It also determines complete safe sets for four explicit maximal-dimension
\(k=2^{40}\) multiplicative rows in the tangent branch and four all-radius
saturated rows \cite[Corollary~1.5 and Proposition~B.2]{ChoShortening}.
Those exact instances do not determine the positive-density identity-profile regime
studied here.
This paper quotes those established frontiers instead of presenting them as
new.  The short endpoint and analytic implications are restated where needed
to make the present compiler auditable.

The object of this companion is the missing classification.  A safe proof
must classify every bad witness, normalize boundary maps by their attained
images, convert each count to the relevant projection, and add the payments
without overlap.  The exact first-match compiler in
\cref{thm:first-match} formalizes these requirements.  The target-exact
profile-closure meta-conjecture \cref{meta:profile-envelope} asks for one
family-uniform, non-oracular atlas schema whose specializations determine
every target comparison.  It is a meta-conjecture about a proposed proof
architecture, not a substitute for a statement about the true numerator.

The stable finite conjectures are therefore stated directly.  Here
\(B_C^{\rm list}(a)\), defined formally in \cref{sec:setup}, denotes the
maximum number of codewords of \(C\) agreeing with one received word on at
least \(a\) coordinates.  The benchmark codes \(C_{\rm KB}\) and
\(C_{\rm M}\), their fields, domains, and targets are defined in
\cref{sec:finite}.  At
\(n=2^{21}\), \(k=2^{20}\), they are
\[
 \begin{aligned}
 B_{C_{\rm KB}}^{\MCA}(1116048)&\le274980728111395087,&
 B_{C_{\rm M}}^{\MCA}(1116024)&\le16777215,\\
 B_{C_{\rm KB}}^{\rm list}(1116047)&\le274980728111395087,&
 B_{C_{\rm M}}^{\rm list}(1116023)&\le16777215.
 \end{aligned}                                      \tag{1.5}\label{eq:intro-benchmarks}
\]
The top line contains the two MCA inequalities; the bottom line contains the
two auxiliary agreeing-codeword inequalities.  Each preceding agreement is rigorously
unsafe, but the four displayed upper inequalities remain open
(\cref{conj:benchmark}).  Their exact margins, periodic audit, random-fiber
heuristic, and falsifiers are recorded in \cref{sec:finite}.

The new theorems concern the interfaces and barriers.  We prove
realized-image moment and support-to-slope incidence inequalities, the
necessary moment orders \(94196,94991,641594,680397\) for the stated
ceiling-normalized route, a general projective incidence theorem, and a
conditional correction-space diagnostic with dimensions near
\(9.1\cdot10^5\).  We also isolate the moving-scale consequence of the
Sidon--Balog--Szemer\'edi--Gowers--Boolean-cube implication already proved in
\cite[Section~10]{ChoShortening}.  Closing any adjacent row still requires
realized-image primitive-fiber control, residual slope or codeword payments,
periodic/descent/extension/rank routing, and exact first-match add-back.

The pole-adjusted identity profile has a precise conjectural asymptotic shape.  Write
\(H_2(x)=-x\log_2x-(1-x)\log_2(1-x)\) for binary entropy, and let
\(\rho_n=k_n/n\), \(\beta_n=\log_2\abs{\B_n}\), and
\[
 g^*(r,b)=\sup\bigl(\{g\in[0,1-r]:H_2(r+g)\ge bg\}\cup\{0\}\bigr).
                                                               \tag{1.6}\label{eq:intro-gstar}
\]
For a declared target sequence \(0<\varepsilon_n\le1\), put
\(q_n=\abs{\F_n}\), \(Q_n=q_n-n\), and
\(B_n^*=\floor{\varepsilon_nq_n}\).  On the nonzero pole-admissible branch
\(B_n^*\ge1\) and \(Q_n>k_nB_n^*\), let \(T_{{\rm pole},n}\) be the exact threshold in
\eqref{eq:pole-threshold}, set
\(\tau_n^{\rm pole}=n^{-1}\log_2T_{{\rm pole},n}\), and define
\[
 g_{{\rm pole},n}=\sup\left(
 \left\{g\in[0,1-\rho_n^+]:
 H_2(\rho_n^++g)-\beta_ng\ge\tau_n^{\rm pole}\right\}\cup\{0\}
 \right),\qquad \rho_n^+=\frac{k_n+1}{n}.          \tag{1.7}\label{eq:intro-pole-crossing}
\]
The exact ceiling test uses \(T_{{\rm pole},n}-1\), rather than
\(T_{{\rm pole},n}\); the normalized logarithmic levels differ by at most
\(1/n\).  A matching exact certificate bracket of width \(o(g_n^*n)\) at
this crossing gives
\[
 \delta_{C_n,\mathrm{off,sup}}^*
 =1-\rho_n-g_{{\rm pole},n}+o(g_n^*),\qquad
 g_n^*=g^*(\rho_n,\beta_n).                        \tag{1.8}\label{eq:intro-pole-prediction}
\]
If in addition the raw budget and the pole-conversion factor are
subexponential,
\[
 \frac1n\log_2B_n^*=o(1),\qquad
 \frac1n\log_2\frac{Q_n-k_n}{Q_n-k_nB_n^*}=o(1), \tag{1.9}\label{eq:intro-pole-condition}
\]
then \cref{lem:crossing-stability} gives
\(g_{{\rm pole},n}=g_n^*+o(g_n^*)\).  Thus the sharper prediction becomes
\[
 \delta_{C_n,\mathrm{off,sup}}^*
 =1-\rho_n-g_n^*+o(g_n^*).                        \tag{1.10}\label{eq:intro-sharp-prediction}
\]
At constant rate \(\rho_n=\rho\) this implies the coarser formula
\[
 \boxed{\displaystyle
 \delta_{C_n,\mathrm{off,sup}}^*
 =1-\rho-g^*(\rho,\log_2\abs{\B_n})+o(1)}.
                                                               \tag{1.11}\label{eq:intro-coarse-prediction}
\]
The stronger collision-nonbinding condition
\(k_nB_n^*/Q_n=o(1)\) implies the second condition in
\eqref{eq:intro-pole-condition}.  Raw subexponential budget alone does not:
there is no off-domain pole when \(Q_n=0\), and the certified pole lower
bound never exceeds the budget when \(0<Q_n\le k_nB_n^*\), regardless of
the locator-list size.

These conclusions do not require every multiplicative loss to be
\(2^{o(g_n^*n)}\).  Near the right crossing, one more agreement changes the
base-two logarithm of the identity mean by \(-\beta_n+O(1)\); consequently a
uniform aggregate factor \(2^{o(n)}\) already moves the crossing by
\(o(n/\beta_n)=o(g_n^*n)\).  Moving-scale control is a useful stronger
statement, especially for finite ledgers, but is not necessary for
\eqref{eq:intro-sharp-prediction}.

Here the tangent branch is \emph{nonbinding} when the previously proved
linear-budget tangent edge has been tested and is not the active constraint
in the crossing window.  To formulate a lower-side hypothesis, fix uniformly
from the public row-family data, before evaluating the unknown MCA numerator
or the eventual safe/unsafe outcome, the non-oracular finite catalogue
\(\mathcal L^{\ne{\rm id}}\) containing every explicit nonidentity lower
construction declared for that family in this manuscript or its pinned
public row specification.  The catalogue is part of the fixed family data;
no subcatalogue may be selected when testing candidate status.
For \(\lambda\in\mathcal L^{\ne{\rm id}}\), write
\(P_{\lambda,n}(a)\le B_{C_n}^{\MCA}(a)\) for its certified slope count and
put
\[
 a_{\lambda,n}^{\rm low}
 =\max\bigl(\{a\in\{k_n+1,\ldots,n\}:
                   P_{\lambda,n}(a)>B_n^*\}\cup\{k_n\}\bigr). \tag{1.12}\label{eq:nonidentity-lower-edge}
\]
Define \(a_{{\rm id},n}^{\rm low}\) by the same formula using the exact
pole-converted identity lower bound \(P_{{\rm id},n}\) from
\eqref{eq:identity-lower}.  A pole-admissible sequence is
\emph{identity-candidate}, relative to this fixed catalogue, when the
empty, saturation, and tangent branches do not bind,
\(a_{{\rm id},n}^{\rm low}\ge k_n+1\), and
\[
 \max_{\lambda\in\mathcal L^{\ne{\rm id}}}
 a_{\lambda,n}^{\rm low}\le a_{{\rm id},n}^{\rm low}           \tag{1.13}\label{eq:identity-candidate-branch}
\]
for all sufficiently large \(n\); the maximum is \(k_n\) if the catalogue
is empty.  This definition uses only realized lower constructions and does
not assume an upper envelope.
\begin{conjecture}[Identity dominance on candidate smooth branches]
\label{conj:identity-frontier}
Let \(C_n=\RS_{\F_n}(D_n,k_n)\) be an official full-field, positive-density
admissibly smooth multiplicative or Chebyshev twin-coset sequence of fixed
rate, with \(B_n^*\ge1\).  Suppose the rows are nonsaturated, the tangent
branch is nonbinding, \(Q_n>k_nB_n^*\), and
\eqref{eq:intro-pole-condition} holds, and the sequence is
identity-candidate relative to the fixed declared catalogue above.  Then there
are integer sequences
\(a_{-,n}<a_{+,n}\) in \(\{k_n+1,\ldots,n\}\) and a family-uniform,
non-oracular profile atlas whose specialization at \(a_{+,n}\) is
witness-exhaustive and has exact first-match payments summing to at most
\(B_n^*\).  Together with the pole-converted identity lower certificate at
\(a_{-,n}\), these give
\[
 B_{C_n}^{\MCA}(a_{-,n})>B_n^*,\qquad
 B_{C_n}^{\MCA}(a_{+,n})\le B_n^*,                 \tag{1.14}\label{eq:intro-exact-bracket}
\]
and
\[
 a_{\pm,n}=k_n+1+g_{{\rm pole},n}n+o(g_n^*n)
 =k_n+1+g_n^*n+o(g_n^*n).                         \tag{1.15}\label{eq:intro-bracket-location}
\]
Consequently the first safe agreement is
\(k_n+1+g_n^*n+o(g_n^*n)\), and the official supremum satisfies
\eqref{eq:intro-sharp-prediction}.  The same assertion applies to a circle
presentation only through a proved agreement-preserving monomial equivalence.
\end{conjecture}
We call a sequence satisfying this conclusion \emph{identity-dominant}.
Thus identity dominance is the predicted conclusion, not part of the
lower-side hypothesis.  Section~\ref{sec:asymptotic} proves the numerical
conclusion from a matching exact certificate bracket; the conjectural
content is witness-exhaustive closure with the asserted payments and bracket.

These are conditional identity-candidate predictions, not unconditional
frontiers: quotient, partial-occupancy, planted, image-collapse, tangent,
or ray-saturation profiles may move the crossing.  Thus the general predicted
answer is the right crossing of the complete image-normalized profile
envelope, after separately testing the empty-budget, tangent, and saturation
branches.

No theorem here improves the safe threshold in \cite{ChoShortening}.  What is
proved here is the compiler and barrier calculus, including the conditional
moving-scale implication and the pole-aware crossing reduction; what remains open
is the Sidon payment, the residual projection compiler, the other algebraic
payments, the four adjacent inequalities, and uniform smooth/circle closure.
The preliminary Prize welcomes partial progress, but any claimed resolution
must survive scientific peer review \cite{ProximityPrize}.

\section{Exact certificate compiler}\label{sec:setup}

Retain the fields, code, challenge set, and budget of \cref{sec:intro}.  A
word \(u\in\F^D\) is \emph{explained} on \(S\subseteq D\) if
\(u|_S=f|_S\) for some \(f\in\F[X]_{<k}\); a pair is simultaneously
explained if both words are, possibly by different polynomials.

For an integer agreement \(a\), an MCA witness for a received pair
\(r=(r_0,r_1)\) is a triple \(w=(\gamma,S,h)\) such that
\[
 \gamma\in\Gam,\qquad \abs S\ge a,\qquad
 (r_0+\gamma r_1)|_S=h|_S,\quad \deg h<k,
\]
and \((r_0,r_1)\) is not simultaneously explained on \(S\).  Let
\(\mathcal W_r(a)\) be the set of witnesses and let
\(Z_r(a)=\{\gamma:(\gamma,S,h)\in\mathcal W_r(a)\}\).  Thus
\[
 B_{C,\Gam}^{\MCA}(a)=\max_{r\in(\F^D)^2}\abs{Z_r(a)}.
\]

The following endpoint conversion is
\cite[Theorem~1.4, Remark~4.4 and Lemma~B.1]{ChoShortening}, restated because
a subcapacity formula alone loses the all-radius branch.

\begin{proposition}[Exact staircase conversion]\label{prop:staircase}
Let \(C\) be a proper \([n,k]\) MDS code.  Then
\(a\mapsto B_{C,\Gam}^{\MCA}(a)\) is nonincreasing,
\[
 B_{C,\Gam}^{\MCA}(a)=B_{C,\Gam}^{\MCA}(k+1)
 \quad(0\le a\le k+1),
 \qquad B_{C,\Gam}^{\MCA}(n)=1.                    \tag{2.1}\label{eq:endpoints}
\]
If \(B^*=0\), no radius is safe.  If
\(B_{C,\Gam}^{\MCA}(k+1)\le B^*\), the complete official safe set is
\([0,1]\).  Otherwise \(B^*\ge1\), and the first safe agreement
\[
 a^*=\min\{a\in\{k+1,\ldots,n\}:B_{C,\Gam}^{\MCA}(a)\le B^*\}
\]
is well defined.  Its complete real safe set is
\[
 \left[0,\frac{n-a^*+1}{n}\right),                 \tag{2.2}\label{eq:safe-set-conversion}
\]
whose right endpoint is unsafe; the largest safe grid radius is
\((n-a^*)/n\).
\end{proposition}

\begin{proof}
This is the cited endpoint theorem and its exact integer-to-radius conversion.
\end{proof}

For a support \(S\), write
\(Q_S(X)=\prod_{x\in S}(X-x)\).  If \(m=\abs S\) and \(K\le m\), the first
\(w=m-K\) nonleading coefficients define the locator-prefix map
\[
 \Phi_w(S)=(c_1(S),\ldots,c_w(S))\in\B^w,
 \qquad Q_S=X^m+c_1(S)X^{m-1}+\cdots+c_m(S).
\]
Pigeonholing its fibers gives a received word with at least
\(\ceil{\binom nm\abs\B^{-w}}\) agreeing codewords of dimension \(K\).
For \(K=k+1\), a collision-aware simple-pole construction converts the list
to MCA slopes.  In a full-field row, with \(q=\abs\F>n\), its exact
consequence at agreement \(a\) is
\[
 L_{\rm id}(a)=\ceil{\binom na\abs{\B}^{-(a-k-1)}},\qquad
 P_{\rm id}(a)=
 \ceil{\frac{L_{\rm id}(a)(q-n)}{q-n+k(L_{\rm id}(a)-1)}}
 \le B_C^{\MCA}(a).                                \tag{2.3}\label{eq:identity-lower}
\]
These exact lower mechanisms and their nontriviality checks
are established in \cite[Proposition~8.1 and Theorem~8.2]{ChoShortening}.  We use their
certified unsafe agreements below, but never interpret the next, smaller
lower value as a safe bound.  In particular, the first agreement at which
\(P_{\rm id}(a)\le B^*\) is only an identity-route candidate; identifying it
with the first safe agreement requires a closed upper atlas.

The following elementary conversion is essential when the challenge field
is not so large that pole collisions are negligible.

\begin{lemma}[Exact pole conversion]\label{lem:exact-pole-conversion}
Let \(B\ge1\), \(Q=q-n>0\), and
\[
 M_{Q,k}(L)=\ceil{\frac{LQ}{Q+k(L-1)}}\qquad(L\in\mathbb Z_{\ge1}).
\]
Then
\[
 M_{Q,k}(L)>B
 \quad\Longleftrightarrow\quad
 L(Q-kB)>B(Q-k).                                  \tag{2.4}\label{eq:pole-equivalence}
\]
If \(Q\le kB\), then \(M_{Q,k}(L)\le B\) for every \(L\ge1\).  If
\(Q>kB\), the least integer \(L\) for which \(M_{Q,k}(L)>B\) is
\[
 T_{\rm pole}(B;Q,k)
 =\floor{\frac{B(Q-k)}{Q-kB}}+1.                 \tag{2.5}\label{eq:pole-threshold}
\]
Consequently the certified pole lower bound in \eqref{eq:identity-lower}
exceeds \(B\) exactly when
\[
 L_{\rm id}(a)\ge T_{\rm pole}(B;Q,k)
 \quad\Longleftrightarrow\quad
 \binom na\abs{\B}^{-(a-k-1)}>T_{\rm pole}(B;Q,k)-1. \tag{2.6}\label{eq:exact-pole-mean}
\]
\end{lemma}

\begin{proof}
Since \(B\) is integral, \(\ceil{x}>B\) if and only if \(x>B\); clearing
the positive denominator gives \eqref{eq:pole-equivalence}.  If
\(Q\le kB\), then for \(L\ge1\),
\(L(Q-kB)\le Q-kB\le B(Q-k)\), so the strict inequality cannot hold.
If \(Q>kB\), then also \(Q>k\), and division by \(Q-kB\) gives
\eqref{eq:pole-threshold}.  The last equivalence uses
\(L_{\rm id}(a)=\ceil{\binom na\abs{\B}^{-(a-k-1)}}\).
\end{proof}

\begin{definition}[Admissibly smooth row]\label{def:admissibly-smooth}
An RS row is \emph{admissibly smooth for a declared profile schema} if its
public data include a finite specified family of polynomial folding maps
\(\pi:D\to D_\pi\), each complete and \(c_\pi\)-to-one in the sense fixed
in \cref{sec:intro}, with \(c_\pi\), every scale, and every occupancy
parameter used by the schema included in those data.
No folding map may be chosen after a heavy fiber or bad received line is
observed.
\end{definition}

\begin{definition}[Declared profile rule]\label{def:profile-rule}
Fix public row data \(P=(\B\subseteq\F,D,k,\Gam,a,B^*)\).  A declared profile
rule is a finite ordered list \(\mathcal A(P)=(\mathcal P_1,\ldots,\mathcal
P_s)\) of predicates on pairs \((r,w)\), where \(r\in(\F^D)^2\) and
\(w\in\mathcal W_r(a)\).  The predicates and their order are specified from
the public data before the values \(\abs{Z_r(a)}\), a maximizing received
line, or the truth of the target inequality are known.  An instantiation is
\emph{auditable} if every predicate is given by explicit algebraic equalities,
inequalities, rank conditions, field-of-definition conditions, or occupancy
conditions involving the received line, the support locator, and the declared
folding maps on \(D\).
\end{definition}

A concrete atlas must display each cell's parameters, predicate, order,
projection to affine slopes, and proved integer payment.  Defining a final
cell as a residual may establish set-theoretic coverage, but it does not pay
that residual.  Singleton predicates chosen from the computed bad set, or
predicates involving the unknown numerator or target outcome, are excluded.
Thus the theorem below is only the exact compiler; exhaustion and payment are
the substantive work.

For a rule and received line, let
\(\mathcal C_i(r)=\{w\in\mathcal W_r(a):\mathcal P_i(r,w)\}\) and let
\(Z_i(r)\) be its projection to \(\Gam\).  The rule is
\emph{witness-exhaustive} for \(r\) if
\(\mathcal W_r(a)=\bigcup_i\mathcal C_i(r)\).  Its first-match slope cells are
\[
 Z_i^\circ(r)=Z_i(r)\setminus\bigcup_{j<i}Z_j(r).
\]
A \emph{payment} is a proved integer inequality
\(\abs{Z_i^\circ(r)}\le U_i(r)\), uniform when it holds for every received
line.  The profile-closure meta-conjecture uses the stronger notion of an
\emph{admissible algebraic payment}.  Such a payment is a specialization of a
row-family-uniform formula
\[
 \abs{Z_i^\circ(r)}\le
 F_i\bigl(P;\Xi_{i,1}(r),\ldots,\Xi_{i,t_i}(r)\bigr),
\]
where the integer formula \(F_i\), the source-side statistics \(\Xi_{i,j}\),
and a proof of the inequality are declared before specialization to a
maximizing received line.  Permitted statistics are explicitly named
algebraic data: parameter-space sizes, degrees, ranks, field-orbit sizes,
folding occupancies, attained images and certified fiber moments of
predeclared maps, and proved incidence multiplicities.  An attained-image or
fiber statistic may enter only through a separately stated row-uniform
estimate.

Neither \(F_i\) nor any \(\Xi_{i,j}\) may invoke, recover, or enumerate
\(Z_i^\circ(r)\), \(Z_r(a)\), their cardinalities, an analogous list-codeword
projection, the true numerator, a maximizing received datum, or the truth of
the target comparison.  Optimization or exhaustive enumeration over received
data, witnesses, slopes, or agreeing codewords is likewise excluded.  Thus
the one-cell choice \(U_1(r)=\abs{Z_r(a)}\), and any disguised equivalent, is
not admissible.  A set-theoretic residual cell is allowed, but it still
requires an independently proved admissible payment.

\begin{theorem}[Exact first-match compiler]\label{thm:first-match}
Suppose a single declared rule \(\mathcal A(P)\), fixed from the public row
data, is witness-exhaustive for every received line and has valid first-match
payments satisfying \(\sum_iU_i(r)\le U\).  Then
\(B_{C,\Gam}^{\MCA}(a)\le U\).  In particular, \(U\le B^*\) is a safe
certificate.

\end{theorem}

\begin{proof}
The sets \(Z_i^\circ(r)\) are disjoint and their union is \(Z_r(a)\).
Therefore \(\abs{Z_r(a)}=\sum_i\abs{Z_i^\circ(r)}\le\sum_iU_i(r)\le U\).
Maximize over \(r\).
\end{proof}

For an auxiliary list row and received word \(u\in\F^D\), define
\[
 \Lambda_C(u,a)=\{c\in C:\abs{\{x\in D:u_x=c_x\}}\ge a\},\qquad
 B_C^{\rm list}(a)=\max_{u\in\F^D}\abs{\Lambda_C(u,a)}.
\]
A declared list rule classifies triples \((u,c,S)\) and projects its cells to
agreeing codewords.  If \(Y_i^\circ(u)\) are the first-match codeword cells
and \(\abs{Y_i^\circ(u)}\le V_i(u)\) with \(\sum_iV_i(u)\le V\), the same
disjoint-union proof gives \(B_C^{\rm list}(a)\le V\).  This projection
change is essential: an MCA ray or slope payment cannot substitute for a
required agreeing-codeword payment, or conversely.  For an atlas at agreement
\(a\), set
\[
\mathcal E_{\mathcal A}(P,a)=\max_r\sum_iU_i(P,r;a), \tag{2.7}\label{eq:profile-envelope-value}
\]
where the maximum ranges over received lines in the MCA case and received
words in the list case.

\begin{metaconjecture}[Target-exact profile-closure program]\label{meta:profile-envelope}
For each declared unbounded-length admissibly smooth multiplicative or
Chebyshev row family
and each public target rule \(P\mapsto B^*(P)\), there is one finite,
auditable algebraic atlas schema fixed uniformly over the family.  For every
proper row \(P\) in the family and every \(k+1\le a\le n\), its
specialization is witness-exhaustive and has admissible algebraic payments
satisfying
\[
 \mathcal E_{\mathcal A}(P,a)\le B^*(P)
 \quad\Longleftrightarrow\quad
 B_{C,\Gam}^{\MCA}(a)\le B^*(P).                  \tag{2.8}\label{eq:profile-envelope}
\]
The schema has finite syntax independent of the row length.  It may depend on
the declared family and target rule, but not on a literal encoding of one
finite row or agreement.  The circle analogue is
asserted only after a proved agreement-preserving monomial equivalence.
\end{metaconjecture}

The atlas must be frozen before the numerator is evaluated and may not encode
the eventual answer.  In particular, one schema, rather than a separately
hard-coded lookup table, must work at every agreement in every row of its
declared family.  Thus
\cref{meta:profile-envelope} is a falsifiable proposal for proof
architecture: an escaping witness family or an unpayable first-match cell
refutes a proposed instantiation.  Independently of this architecture,
\cref{conj:benchmark} consists of four explicit inequalities for the true
MCA or list numerators.

For reference, in a full-field row with \(q=\abs\F>n\) and \(B^*\ge1\), the
identity route singles out
\[
 a_{\rm id}^*
 =\min\{a\in\{k+1,\ldots,n\}:P_{\rm id}(a)\le B^*\}. \tag{2.9}\label{eq:identity-candidate}
\]
When the empty, saturation, and tangent branches do not bind and every lower
edge in the fixed declared nonidentity catalogue is no later than the corresponding
identity edge, this is the identity candidate in the sense of
\eqref{eq:identity-candidate-branch}.  Its preceding agreement is
unconditionally unsafe by \eqref{eq:identity-lower}; safety at
\(a_{\rm id}^*\) still requires a closed upper atlas and remains conjectural
both for the direct rows in \cref{conj:benchmark} and for the asymptotic
families in \cref{conj:identity-frontier}.

The proposed atlas has four payment groups.  Algebraic routing must cover the
endpoint and tangent branches, declared foldings with full and partial
occupancy, quotient, remainder, and planted factors, and all descent,
extension-orbit, rank-loss, sparse, and auxiliary cells.  Here auxiliary
means in particular correlated-agreement and shortened or interleaved-list
residuals; orbit stabilizers, extension degrees, and named rank
multiplicities remain in the payment.  Primitive cells require estimates on
the attained image together with max-fiber and incidence degrees.  Every
residual chart requires a distinct-slope payment for MCA or a distinct
codeword payment for a list row.  Finally, first-match deduplication must add
the integer payments with a reserve below the target.  In particular, for a
declared \(c\)-to-one folding \(\pi:D\to D'\), the public occupancy datum is
\(\nu_\pi(S)=(\abs{S\cap\pi^{-1}(y)})_{y\in D'}\); ``primitive'' means only
that the witness escaped the preceding declared cells, not that its boundary
map is surjective or flat.

For multiplicative rows the declared maps are \(x\mapsto x^c\); for
Chebyshev twin-coset rows they are \(T_c\).  Circle presentations enter only
after a coordinate bijection and nonzero diagonal scaling give a
support-preserving monomial equivalence with an RS row
\cite{HLP24,ChoShortening}.  The calculus itself is
characteristic-independent, but the norm-one-torus presentation used here is
in odd characteristic; in characteristic two we claim only statements for
explicit RS rows and declared algebraic maps.

\section{Primitive-fiber barriers}\label{sec:primitive}

The primitive problem begins only after the earlier algebraic cells have been
declared.  Its normalization must refer to the full reference slice before
first-match deletion.  Let \(\Om^0\) be such a nonempty finite slice, let
\(\Om^\circ\subseteq\Om^0\) be a residual leaf, and let
\(\Phi:\Om^0\to A\) be a boundary map.  Write
\[
 I=\Phi(\Om^0),\qquad M=\abs{\Om^0},\qquad L=\abs I,
 \qquad \barN=\frac ML,
\]
and, for \(y\in I\), put
\(F_y=\Om^\circ\cap\Phi^{-1}(y)\), \(f_y=\abs{F_y}\), and
\(W=\sum_yf_y=\abs{\Om^\circ}\).  Thus \(\barN\) is the mean fiber size of
the reference slice, not the mean over the smaller residual image.  This is
the normalization used by the finite incidence payment and by the
Gowers/cube implication in \cite{ChoShortening}.

For an integer \(r\ge2\), define the image-normalized residual moment
\[
 \Gamma_r=\frac1L\sum_{y\in I}\left(\frac{f_y}{\barN}\right)^r.
                                                               \tag{3.1}\label{eq:residual-moment}
\]
The elementary maximum bound is
\[
 \max_y f_y\le \barN(L\Gamma_r)^{1/r}.              \tag{3.2}\label{eq:moment-max}
\]
Indeed, the largest summand in \eqref{eq:residual-moment} is at most the
whole sum.  This already shows why an ordinary-moment conjecture strong enough
to make \(\log(L\Gamma_r)/r=o(n)\) would itself prove max-fiber flatness;
the Gowers/cube step is needed for a weaker, low-energy input instead.

The exact support-fiber-to-slope interface also needs an incidence theorem.
Let \(Z^\circ\) be the first-match slope projection of the leaf, and let
\(\mathcal P=\{(u,v)\in(\Om^\circ)^2:\Phi(u)=\Phi(v)\}\).  Suppose an
incidence relation \(\mathcal I\subseteq Z^\circ\times\mathcal P\) has degree
at least \(H\ge1\) at every slope and degree at most \(J\) at every ordered
pair.

\begin{proposition}[Exact primitive incidence payment]\label{prop:primitive-payment}
With the preceding notation,
\[
 \abs{Z^\circ}
 \le \floor{\frac{JW}{H}\,\barN(L\Gamma_r)^{1/r}}. \tag{3.3}\label{eq:primitive-payment}
\]
An independent direct slope bound may be intersected with this one.
\end{proposition}

\begin{proof}
Double counting gives
\(H\abs{Z^\circ}\le J\sum_y f_y^2\).  Since
\(\sum_yf_y^2\le W\max_yf_y\), \eqref{eq:moment-max} proves the result;
the left side is an integer.
\end{proof}

The quantities \(H\) and \(J\) are not cosmetic.  A bound for
\(\max_yf_y\), or even for all same-fiber pairs, does not bound slopes until
one proves how many such pairs a slope supplies and how many slopes a pair may
represent.

For the full fiber distribution \(\Om^\circ=\Om^0\), set
\(N_y=\abs{\Phi^{-1}(y)}\), \(\mu_y=N_y/M\), and
\[
 R=L\max_y\mu_y,
 \qquad \Gamma_r^{\rm full}=L^{r-1}\sum_{y\in I}\mu_y^r.
\]
Here \(R\) is the maximum-fiber overhead relative to \(M/L\), and
\(\Gamma_r^{\rm full}=1\) for a uniform distribution.

\begin{proposition}[Exact moment sandwich]\label{prop:moment-sandwich}
For every integer \(r\ge2\),
\[
 \frac{\log_2\Gamma_r^{\rm full}}{r-1}
 \le\log_2R
 \le\frac{\log_2L+\log_2\Gamma_r^{\rm full}}r.     \tag{3.4}\label{eq:moment-sandwich}
\]
Consequently \(R\le2^\Delta\) follows from
\(\log_2L+\log_2\Gamma_r^{\rm full}\le r\Delta\).
\end{proposition}

\begin{proof}
The inequalities
\(\sum_y\mu_y^r\le(\max_y\mu_y)^{r-1}\) and
\((\max_y\mu_y)^r\le\sum_y\mu_y^r\) give
\(\Gamma_r^{\rm full}\le R^{r-1}\) and
\(R^r\le L\Gamma_r^{\rm full}\), respectively.
\end{proof}

The image \(I\) cannot silently be replaced by the formal codomain \(A\).
Let \(A_0=\abs A\), extend both \(f_y\) and \(N_y\) by zero on
\(A\setminus I\), and put \(\barN^{\rm amb}=M/A_0\).  The residual ambient
moment is
\[
 \Gamma_{r,\mathrm{res}}^{\rm amb}
 =\frac1{A_0}\sum_{y\in A}
   \left(\frac{f_y}{\barN^{\rm amb}}\right)^r
 =\left(\frac{A_0}{L}\right)^{r-1}\Gamma_r.       \tag{3.5}\label{eq:ambient-image}
\]
For the full reference distribution, define
\[
 \Gamma_{r,\mathrm{full}}^{\rm amb}
 =\frac1{A_0}\sum_{y\in A}
   \left(\frac{N_y}{\barN^{\rm amb}}\right)^r
 =A_0^{r-1}\sum_{y\in A}\left(\frac{N_y}{M}\right)^r\ge1. \tag{3.6}\label{eq:ambient-full}
\]
The last inequality follows from the power-mean inequality and
\(\sum_yN_y/M=1\).  Thus an ambient residual-moment estimate implies its
image-normalized residual counterpart, but the converse pays the exact
factor \((A_0/L)^{r-1}\).  At exponent scale one may discard this factor
only after proving \(\log(A_0/L)=o(n)\) at the relevant order.

Two distinct failures are encoded here.  First, the reference map itself may
have a small image \(I\), in which case \(M/L\) is much larger than the
formal mean \(M/A_0\).  Second, first-match deletion may concentrate the
residual leaf inside a few fibers even when the reference image is large.
Resetting the denominator to \(\abs{\Phi(\Om^\circ)}\) after this deletion
would change the conjecture and sever the link to the preceding payments.
The reference-slice normalization keeps both losses visible: image collapse
appears through \(A_0/L\), and residual concentration appears in the moment
\(\Gamma_r\).

\begin{proposition}[Ambient moment-order barrier]\label{prop:ambient-order}
Suppose a prefix map has formal codomain \(\B^w\), and normalize its full
distribution over all \(Q=\abs{\B}^w\) formal values, including empty fibers.
Any argument using only the upper half of \cref{prop:moment-sandwich} and
\(\Gamma_{r,\mathrm{full}}^{\rm amb}\ge1\) to force ambient maximum-fiber overhead at most
\(2^\Delta\) must have
\[
 r\ge\ceil{\frac{w\log_2\abs\B}{\Delta}}.          \tag{3.7}\label{eq:order-floor}
\]
\end{proposition}

\begin{proof}
That route requires
\(\log_2Q+\log_2\Gamma_{r,\mathrm{full}}^{\rm amb}\le r\Delta\).
Discarding the nonnegative second term and using
\(\log_2Q=w\log_2\abs\B\) gives the claim.
\end{proof}

We now state the missing analytic input at a general resolution scale.  Let
\(\B_N\) be a finite field, let \(T_N\subseteq\B_N\) consist of \(N\)
distinct points, let \(1\le d_N\le N\), and choose nonzero weights
\(\omega_N(t)\in\B_N^\times\).  The weighted Vandermonde boundary map is
\[
 \Phi_N(x)=\sum_{t\in T_N}x_t\omega_N(t)
 (1,t,\ldots,t^{d_N-1})^{\mathsf T}.               \tag{3.8}\label{eq:weighted-vandermonde}
\]
Fix \(0<\alpha\le1/2\), take
\(\Omega_N^0=\{x\in\{0,1\}^{T_N}:\sum_tx_t=m_N\}\), where
\(m_N/N\in[\alpha,1-\alpha]\), and let
\(\Omega_N^\circ\subseteq\Omega_N^0\) be a residual leaf selected by one
fixed declared rule.  Set
\[
 I_N=\Phi_N(\Omega_N^0),\quad L_N=\abs{I_N},\quad
 \barN_N=\frac{\abs{\Omega_N^0}}{L_N},\quad
 F_{N,y}=\Omega_N^\circ\cap\Phi_N^{-1}(y).
\]
For \(F\subseteq\Z^{T_N}\), its additive energy is
\(E(F)=\abs{\{(x_1,x_2,x_3,x_4)\in F^4:x_1+x_2=x_3+x_4\}}\).
Compute this energy in the torsion-free group \(\Z^{T_N}\), and put
\(\Delta_{N,y}=E(F_{N,y})/\abs{F_{N,y}}^3\), with value zero for an empty
fiber.  Let \(s_N\to\infty\), with \(s_N\le N\).  For \(q\ge2\) and
\(\sigma>0\), define
\[
 \Gamma^{\rm sid}_{N,q,\sigma;s}
 =\frac1{L_N}\sum_{\Delta_{N,y}<e^{-\sigma s_N}}
   \left(\frac{\abs{F_{N,y}}}{\barN_N}\right)^q.   \tag{3.9}\label{eq:sidon-moment}
\]

\begin{problem}[Primitive Sidon payment]\label{prob:sidon}
For primitive locator-prefix leaves in a fixed, predeclared multiplicative
or Chebyshev atlas, prove that there are integer orders \(q_N\ge2\) such
that, for every fixed \(\sigma>0\), uniformly over the declared leaves,
\[
 \frac{\log L_N}{q_N}=o(s_N),\qquad
 \Gamma^{\rm sid}_{N,q_N,\sigma;s}\le e^{o(s_Nq_N)}. \tag{3.10}\label{eq:sidon-payment}
\]
The primary asymptotic scale is \(s_N=N\).  As a stronger localization
problem one may take \(N=n\) and \(s_n=g_n^*n\), under the necessary
assumption \(g_n^*n\to\infty\); this stronger scale is not required for the
\(o(g_n^*)\) frontier once all aggregate certificate losses are uniformly
\(2^{o(n)}\).
\end{problem}

The adjective primitive is an input to this problem, not a conclusion of an
equidistribution theorem.  In particular, the rule defining the residual leaf
must be fixed before a heavy fiber is known, and the moment is normalized by
the image of the full reference slice.

We import the following consequence of polynomial
Balog--Szemer\'edi--Gowers and Boolean-quasicube difference growth from
\cite[Proposition~10.4]{ChoShortening}.

\begin{proposition}[Imported BSG--Boolean-cube consequence]
\label{prop:high-energy}
There is an absolute constant \(C_0\) such that, for every finite set \(T\),
every nonempty \(F\subseteq\{0,1\}^T\), and every \(K\ge1\),
\[
 E(F)\ge\frac{\abs F^3}{K}\quad\Longrightarrow\quad
 \abs F\le K^{C_0}.                                \tag{3.11}\label{eq:high-energy}
\]
Here energy is computed after embedding the Boolean cube in \(\Z^T\).
\end{proposition}

\begin{corollary}[Moving-scale primitive flatness]
\label{cor:gowers-cube}
If \eqref{eq:sidon-payment} holds at a scale \(s_N\to\infty\), then
\[
 \max_{y\in I_N}\abs{F_{N,y}}\le e^{o(s_N)}\barN_N. \tag{3.12}\label{eq:scale-flatness}
\]
For \(s_N=N\), this is \cite[Theorem~10.5]{ChoShortening}; the additional
observation is that the same implication is valid at the displayed moving
scale.
\end{corollary}

\begin{proof}
If the conclusion fails, then along a subsequence there are \(\eta>0\) and
fibers of size at least \(e^{\eta s_N}\barN_N\).  The ordinary moment from
\eqref{eq:residual-moment}, with \(r=q_N\), is therefore at least
\(L_N^{-1}e^{\eta s_Nq_N}\ge e^{3\eta s_Nq_N/4}\) for large \(N\), by
\(\log L_N/q_N=o(s_N)\).  Choose \(0<\sigma<\eta/(4C_0)\).  The low-energy
contribution is at most \(e^{\eta s_Nq_N/4}\), so the high-energy
contribution is at least \(e^{\eta s_Nq_N/2}\).  Some high-energy fiber then
satisfies \(\abs{F_{N,y}}/\barN_N\ge e^{\eta s_N/2}\).  Since
\(I_N=\Phi_N(\Omega_N^0)\), one has \(L_N\le\abs{\Omega_N^0}\) and hence
\(\barN_N\ge1\); this gives \(\abs{F_{N,y}}\ge e^{\eta s_N/2}\).
The same fiber has normalized energy at least \(e^{-\sigma s_N}\), so
\cref{prop:high-energy}, with \(K=e^{\sigma s_N}\), gives
\(\abs{F_{N,y}}\le e^{C_0\sigma s_N}<e^{\eta s_N/4}\), a contradiction.
\end{proof}

\begin{corollary}[Analytic-to-slope interface]\label{cor:analytic-to-slope}
Assume \cref{prob:sidon} and suppose the primitive algebraic analysis proves
an incidence relation with parameters \(H_N,J_N\ge1\).  If
\(W_N=\abs{\Omega_N^\circ}\) and \(Z_N^\circ\) is the realized slope set,
then
\[
 \abs{Z_N^\circ}\le
 \frac{J_NW_N}{H_N}\,e^{o(s_N)}\barN_N.           \tag{3.13}\label{eq:analytic-slope}
\]
\end{corollary}

\begin{proof}
Use the incidence double count from \cref{prop:primitive-payment} and
\cref{cor:gowers-cube} in place of \eqref{eq:moment-max}.
\end{proof}

This implication does not solve \cref{prob:sidon}, prove reference-image
coverage, supply the incidence degrees, or close the profile ledger.  At the
coarse scale it yields only an \(e^{o(N)}\) loss, which does not fit a finite
margin of a few bits, but would be asymptotically sufficient for the sharp
identity crossing if every other loss were also uniform and \(e^{o(n)}\).
For the optional stronger specialization \(s_n=g_n^*n\), the energy cutoff
itself must be \(e^{-\sigma g_n^*n}\); changing only the moment order while
retaining the coarse cutoff is not a valid moving-scale argument.  The phrase
``Gowers/cube'' here means BSG followed by Boolean-quasicube difference
growth, not a Gowers-uniformity norm or cube-counting theorem.


\section{Ray-projection barriers}\label{sec:rays}

A same-fiber pair is often called a shift pair.  Max-fiber control gives the
support-level estimate
\(\sum_y\binom{N_y}{2}\le N_{\max}M/2\), and hence at most
\(\kappa M^2/(2L)\) when \(N_{\max}\le\kappa M/L\).  This does not supply
the ray incidence required in \cref{prop:primitive-payment}.

The next incidence bound identifies exactly what a one-parameter
moving-root proof uses.  It is stated for a general-dimensional family to
show how quickly that argument deteriorates.

\begin{theorem}[Projective moving-root incidence]\label{thm:projective-incidence}
Let \(d\ge1\), let \(V\le\F_q[X]\) be a \(d\)-dimensional polynomial subspace, and let
\(G\subseteq D\) be its common zero set, of size \(g\).  Let \(h\ge1\), and suppose
\(\mathcal A\subseteq\mathbb P(V)\) and every \([f]\in\mathcal A\) has at
least \(h\) zeros in \(D\setminus G\).  Then
\[
 \abs{\mathcal A}h\le(n-g)\frac{q^{d-1}-1}{q-1}.    \tag{4.1}\label{eq:projective-incidence}
\]
For the affine analogue, let \(\mathcal F=f_0+V\) be an affine polynomial
space with \(\dim V=d\), let \(G\) be the common zero set in \(D\) of every
member of \(\mathcal F\), of size \(g\), and let
\(\mathcal A\subseteq\mathcal F\).  If
every member of \(\mathcal A\) has at least \(h\) zeros in \(D\setminus G\),
then \(\abs{\mathcal A}h\le(n-g)q^{d-1}\).
\end{theorem}

\begin{proof}
Count pairs \(([f],x)\) with \([f]\in\mathcal A\),
\(x\in D\setminus G\), and \(f(x)=0\).  Every projective member contributes
at least \(h\) pairs.  At a fixed \(x\notin G\), evaluation is a nonzero
linear functional on \(V\), so its projective kernel has
\((q^{d-1}-1)/(q-1)\) points.  This proves \eqref{eq:projective-incidence}.
For the affine statement, evaluation at a point outside \(G\) is either a
nonconstant affine linear functional, whose zero set has \(q^{d-1}\) points,
or a nonzero constant, whose zero set is empty.  The same double count applies.
\end{proof}

\begin{corollary}[Moving-root pencil]\label{cor:moving-root}
Let \(A,B\in\F_q[X]\) be linearly independent, let \(G\subseteq D\) be their
common zero set, of size \(g\), and consider
\(L_{[s:t]}=sA+tB\).  If every counted parameter has at least \(h\) zeros
in \(D\setminus G\), then there are at most \(\floor{(n-g)/h}\) counted
parameters.  If every counted member is squarefree, split over \(D\), and
has degree \(j>g\), the bound is \(\floor{(n-g)/(j-g)}\).
\end{corollary}

\begin{proof}
Use \cref{thm:projective-incidence} with \(d=2\).
\end{proof}

The pencil bound is strong because a point outside the base locus chooses at
most one projective parameter.  In dimension \(d\), it instead chooses a
projective hyperplane, and the factor in
\eqref{eq:projective-incidence} is its exact one-coordinate worst-case
multiplicity.  Thus for large \(d\) the raw bound is essentially
\((n-g)q^{d-2}/h\).  A useful high-dimensional compiler must exploit several
roots, restricted factorization or rank, or a lower-dimensional
parametrization of the rays actually realized by MCA witnesses.

\begin{proposition}[Dimension diagnostic for the declared correction spaces]\label{prop:not-pencil}
At the four candidate adjacent agreements of \cref{tab:benchmark}, assume
the declared prefix reduction parametrizes a residual chart by the full
affine correction space
\(A\in\F[X]_{\le\omega-w-1}\), where \(\omega=n-a^*\) and
\(w=a^*-K\).  Hence its affine dimension is
\(\omega-w=n-2a^*+K\), equal respectively to
\[
 913633,\qquad913634,\qquad913681,\qquad913682.       \tag{4.2}\label{eq:benchmark-dimensions}
\]
Consequently the pencil case of \cref{thm:projective-incidence} cannot pay
this declared full correction space without a further structural reduction.
\end{proposition}

\begin{proof}
Under the stated parametrization there is one free coefficient in each degree from zero through
\(\omega-w-1\).  Substitution of \(n=2097152\), the agreements and the values
\(K=k+1\) for MCA and \(K=k\) for the auxiliary list rows gives the four
integers.  A family of these dimensions is not a two-dimensional projective
subspace or a one-dimensional affine pencil.
\end{proof}

The proposition verifies the correction-space arithmetic conditional on the
declared parametrization.  The release source package contains
\path{verify_profile_barriers.py}, which independently recomputes
\eqref{eq:benchmark-dimensions}.  The lower-reserve artifacts in
\cite{RSMCARepo} record the separate unsafe-side inputs; they are not evidence
for witness exhaustion.  Proving that the declared charts exhaust the
residual witnesses is part of \cref{prob:finite-closure}, not a consequence
of the dimension count.

\begin{problem}[Balanced-core ray payment]\label{prob:balanced-core}
Every residual balanced-core cell in the benchmark rows admits an explicit
witness-exhaustive decomposition into structured charts such that each chart
is either a projective pencil covered by \cref{cor:moving-root}, a controlled
finite union of such pencils, or a higher-dimensional chart with a proved
distinct-slope payment.  After first-match deduplication and exact incidence
multiplicities, the total balanced-core payment fits the row allocation in
\cref{prob:finite-closure}.
\end{problem}

This is a slope statement, not merely a factorization statement.  It is
falsified by a single residual chart whose actual slope projection exceeds
the allocated budget and which cannot be routed to an earlier paid cell.

\section{Finite benchmark rows}\label{sec:finite}

We now record four adjacent rows that make the preceding obstructions
quantitative.  Set
\[
 n=2^{21}=2097152,\qquad k=2^{20}=1048576,
\]
and define
\[
 p_{\rm KB}=2^{31}-2^{24}+1=2130706433,
 \qquad p_{\rm M}=2^{31}-1=2147483647.
\]
Let \(D_{\rm KB}\subseteq\F_{p_{\rm KB}}^\times\) be the multiplicative
subgroup of order \(n\), and put
\(C_{\rm KB}=\RS_{\F_{p_{\rm KB}^6}}(D_{\rm KB},k)\).  Let \(D_{\rm M}\)
be the Chebyshev twin-coset domain of size \(n\) specified in
\cite{ChoShortening,RSMCARepo}, and put
\(C_{\rm M}=\RS_{\F_{p_{\rm M}^4}}(D_{\rm M},k)\).  Their integer budgets are
\[
 B_{\rm KB}^*=\floor{p_{\rm KB}^6/2^{128}}
 =274980728111395087,\qquad
 B_{\rm M}^*=\floor{p_{\rm M}^4/2^{100}}=16777215. \tag{5.1}\label{eq:budgets}
\]
The Mersenne row is an auxiliary \(2^{-100}\) comparison; at target
\(2^{-128}\) this challenge field has zero integer budget.

The list numerator and its separate first-match codeword projection were
defined in \cref{sec:setup}.  The two list rows below are auxiliary inputs,
not Proximity Prize MCA thresholds.

\begin{table}[t]
\centering
\small
\caption{Proved unsafe agreements, adjacent candidate agreements, and their
candidate threshold intervals.  Only the first and third rows are MCA
rows.}\label{tab:benchmark}
\begin{tabular}{@{}lrrl@{}}
\toprule
row and target & proved unsafe \(a_0\) & candidate \(a^*=a_0+1\)
& candidate threshold interval\\
\midrule
KoalaBear MCA, \(2^{-128}\) &1116047&1116048&
\([0,981105/2097152)\)\\
KoalaBear list, \(2^{-128}\)&1116046&1116047&
\([0,490553/1048576)\)\\
Mersenne--31 MCA, \(2^{-100}\)&1116023&1116024&
\([0,981129/2097152)\)\\
Mersenne--31 list, \(2^{-100}\)&1116022&1116023&
\([0,490565/1048576)\)\\
\bottomrule
\end{tabular}
\end{table}

The word ``candidate'' is essential.  The exact prefix construction proves
that each \(a_0\) is unsafe, and hence that every smaller agreement is unsafe.
At \(a^*\), that particular lower route drops below the budget.  This is not
an upper bound for other witnesses.

For completeness, the exact lower-route numerators at \(a_0\), in the order
of \cref{tab:benchmark}, are
\[
 \begin{aligned}
  &138634741058327852652,\qquad157702518233425975347,\\
  &4281388998575706,\qquad4870025984688527.
 \end{aligned}                                      \tag{5.2}\label{eq:unsafe-lower-values}
\]
Each is strictly larger than its row budget in \eqref{eq:budgets}.  The MCA
values use the dimension-\(k+1\) locator-prefix list and the simple-pole
conversion; the list values use dimension \(k\) directly.  In the two MCA
rows the pole-separation condition is satisfied, so no collision factor is
lost at these displayed agreements \cite[Proposition~8.6]{ChoShortening};
the auxiliary values are recorded in \cite{RSMCARepo}, and
\path{verify_profile_barriers.py} in the release source package independently
recomputes all four displayed integers.

The location of the one-step edge has a simple exact explanation.  For the
unrounded prefix mean, put
\(\mu_K(a)=\binom na p^{-(a-K)}\), with \(K=k+1\) in the MCA route and
\(K=k\) in the auxiliary list route.  Then
\(\mu_K(a)/\mu_K(a+1)=p(a+1)/(n-a)\).  In the four rate-one-half rows this
is a drop of roughly \(31\) bits in one agreement step.  The identity route
therefore identifies the microscopic candidate scale, but the ratio supplies
no safe-side bound and does not rule out a later nonidentity crossing.

\begin{conjecture}[Benchmark adjacent safe inequalities]\label{conj:benchmark}
For the two MCA rows,
\[
 B_{C_{\rm KB}}^{\MCA}(1116048)\le274980728111395087,
 \qquad
 B_{C_{\rm M}}^{\MCA}(1116024)\le16777215.           \tag{5.3}\label{eq:mca-adjacent}
\]
For the two auxiliary list rows,
\[
 B_{C_{\rm KB}}^{\rm list}(1116047)\le274980728111395087,
 \qquad
 B_{C_{\rm M}}^{\rm list}(1116023)\le16777215.      \tag{5.4}\label{eq:list-adjacent}
\]
\end{conjecture}

These are direct, falsifiable statements; they do not define their own proof
method.  If either MCA inequality is proved, the corresponding preceding
unsafe certificate and monotonicity make the displayed \(a^*\) the first safe
agreement.  The exact integer-to-radius conversion then gives the associated
half-open interval in \cref{tab:benchmark}.  The same observation applies to
the auxiliary list staircase.

To fix the notation used in the diagnostics, put
\[
 L_K(a)=\ceil{\binom na p^{-(a-K)}},\qquad
 P_{\rm cand}(a)=
 \begin{cases}
 \displaystyle\ceil{\frac{L_{k+1}(a)(q-n)}
 {q-n+k(L_{k+1}(a)-1)}},&\text{MCA},\\[6pt]
 L_k(a),&\text{auxiliary list},
 \end{cases}
\]
using the base prime \(p\) and challenge-field order \(q\) of the row.  With
\(K=k+1\) for MCA and \(K=k\) for list, let
\[
 \Delta_{\rm ceil}=\log_2\frac{B^*}{P_{\rm cand}(a^*)},\qquad
 w=a^*-K,\qquad
 H_{\rm ext}=\floor{\frac{p^wB^*}{\binom n{a^*}}}.
\]
The candidate values and remaining margins are as follows.  They are
reproduced by \path{verify_profile_barriers.py} in the release source package.

\begin{table}[t]
\centering
\small
\caption{Candidate-agreement diagnostics.  The necessary order refers only
to the ceiling-normalized ambient-moment route of
\cref{prop:ambient-order}.}\label{tab:diagnostics}
\begin{tabular}{@{}lrrrr@{}}
\toprule
row & \(P_{\rm cand}(a^*)\)& margin bits & necessary order & extension headroom\\
\midrule
KoalaBear MCA &57198030366&22.1969&94196&4807520\\
KoalaBear list&65065153468&22.0109&94991&4226236\\
Mersenne--31 MCA&1752700&3.2589&641594&9\\
Mersenne--31 list&1993678&3.0730&680397&8\\
\bottomrule
\end{tabular}
\end{table}

The periodic lower-family audit has deliberately limited coverage.  In the
repository labels, \(\mathrm{Gfloor}\) and \(\mathrm{Gceil}\) use
\(m=\floor{a/c}\) and \(m=\ceil{a/c}\), respectively; \(\mathrm{Rem}\)
adds the uncovered remainder, and \(\mathrm{Plant}\) tests a planted quotient
core.  Gate G6 of the exact verifier scans all twenty-one dyadic scales
\(c=2^j\), \(0\le j\le20\), and these four lower families at the historical
packet pairs: \((1116043,1116044)\) for KoalaBear MCA,
\((1116021,1116022)\) for Mersenne--31 MCA, and the unchanged list pairs in
\cref{tab:benchmark}.  Gate G7 separately recomputes the moved MCA identity
pairs in that table and the named Mersenne--31 MCA
\(\mathrm{Gceil}\) watch item at \(c=2048\), which covers agreement
\(1116160\) and has exact mass \(12769758<16777215=B^*\), a deficit of
\(0.3938\) bits.  It does not perform a full twenty-one-scale G6 rescan at
both moved MCA pairs and does not certify safety.  The precise executable is
\path{experimental/scripts/verify_frontier_adjacent_v13_rows.py} at pinned
repository commit \texttt{e190193cebce}.  The release source package contains
a verbatim copy at that path and the four input packets under
\path{experimental/data/certificates/frontier-adjacent/} \cite{RSMCARepo}.
These checks veto only the lower
constructions they actually enumerate; they are neither an upper payment nor
an exhaustion theorem.

There is also a useful but nonrigorous scale check.  If \(M\) objects were
allocated independently and uniformly among \(L\) realized fibers, with
mean \(\mu=M/L\), the Poisson extreme-value factor \(x\) would be predicted
by
\[
 \mu(x\log x-x+1)\approx\log L.
\]
Under the additional full-image assumption, the four rows in table order
give predicted maximum-to-mean factors approximately
\(1.0071,1.0067,2.54,2.43\), while their target-to-mean factors are
approximately
\(4.81\cdot10^6,4.23\cdot10^6,9.57,8.42\).
Thus the random-fiber model favors all four primitive inequalities, most
strongly in the KoalaBear rows.  It is not evidence against effective-image
collapse or algebraic concentration and is never used in a safe-side proof.

The remaining safe-side obligations are not interchangeable.  Algebraic
routing must cover the dense tangent/common-line cell, the full folded and
periodic union, and all field, rank, extension, sparse, correlated-agreement,
shortened, and interleaved-list cells, with orbit stabilizers and every
degree and rank-loss multiplicity retained.
Primitive leaves require a finite image-normalized estimate and the exact
incidence degrees \(H,J\).  Every residual chart then needs a slope payment
for MCA or a codeword payment for a list row.  Finally, the compiler must
deduplicate and add all payments inside the exact integer reserve.

Here the ceiling-normalized margin is the quantity \(\Delta_{\rm ceil}\)
defined above.  A factor \(n^C\) costs
\(21C\) bits, so an unspecified polynomial loss cannot decide any Mersenne
adjacent row and can consume the entire KoalaBear margin even for modest
\(C\).  Applying \cref{prop:ambient-order} at prefix depths
\(67471,67471,67447,67447\) gives the four necessary orders in
\cref{tab:diagnostics}.  This diagnoses only the ceiling-normalized ambient
route; it is not a lower bound on every max-fiber theorem, and a smaller
attained image creates an additional loss.

The extension column has another exact interpretation.  The descended
extension charts in the current ledger have prefix-normalized cost
\(\Delta_0p^e\), where \(p\) is the base-field order, \(e\) is the chart
dimension, and \(\Delta_0\ge1\) is a degree factor.  For each row its
displayed headroom is \(H_{\rm ext}\) as defined above.  The largest displayed
headroom is below \(2^{23}\), while both base primes exceed \(2^{30}\).
Thus no positive-dimensional chart of this form fits: necessarily \(e=0\),
and then \(\Delta_0\) must not exceed the corresponding headroom.  This is a
useful veto, not a payment for all extension witnesses.

\begin{problem}[Finite row-sharp closure]\label{prob:finite-closure}
For either MCA row of \cref{tab:benchmark}, display a witness-exhaustive
declared atlas at the candidate \(a^*\) and prove payments whose exact sum is
at most \(B^*\).  In particular, the proof must:
\begin{enumerate}[label=\textup{(\alph*)}]
\item pay or witness-preservingly route the dense tangent/common-line cell;
the full quotient, remainder, planted, periodic, and partial-occupancy union;
every field-descent and extension orbit with its stabilizer and degree; every
rank-loss multiplicity; and all sparse, correlated-agreement, shortened, and
interleaved-list residuals;
\item prove realized-image primitive-fiber estimates and the incidence
degrees in \cref{prop:primitive-payment};
\item prove a residual compiler for distinct affine slopes or projective rays
in MCA, including the balanced-core payment of \cref{prob:balanced-core};
\item deduplicate by first match and add every integer payment, rounding,
degree, and multiplicity with a reserve that keeps the total at most \(B^*\).
\end{enumerate}
For either auxiliary list row, the corresponding closure problem uses
first-match codeword projections throughout: the balanced core requires a
direct codeword payment, and the ray-incidence item is inapplicable.
\end{problem}

A counterexample to \eqref{eq:mca-adjacent} is a received line with more than
the displayed number of bad slopes; a counterexample to
\eqref{eq:list-adjacent} is a received word with more than the displayed
number of agreeing codewords.  A counterexample to a proposed profile
proof may be weaker: one unpaid first-match cell, a collapsed realized image,
or an underestimated pair-to-slope multiplicity is enough to invalidate that
proof without deciding the true inequality.  Conversely, proving only
\cref{prob:sidon} and \cref{prob:balanced-core} does not close a row; all
four groups in \cref{prob:finite-closure} remain part of the certificate.

\section{Target-aware asymptotic criterion}\label{sec:asymptotic}

The asymptotic formula is best understood as a crossing calculation that a
future certificate must realize on both sides.  Let
\(C_n=\RS_{\F_n}(D_n,k_n)\), where \(D_n\subseteq\B_n\subseteq\F_n\), and
assume \(k_n/n\to\rho\in(0,1)\).  Throughout this section the challenge is
the official full field, \(\Gam_n=\F_n\), and \(\B_n\) is the designated
generated coefficient-and-folding field fixed in \cref{sec:intro}; a sampled
challenge or an arbitrary enlargement of \(\B_n\) is not covered by the
formula below.  On the nonzero-budget branch, put
\[
 \begin{gathered}
 B_n(a)=B_{C_n}^{\MCA}(a),\qquad
 \rho_n=\frac{k_n}{n},\qquad \rho_n^+=\frac{k_n+1}{n},\qquad
 \beta_n=\log_2\abs{\B_n},\\
 q_n=\abs{\F_n},\qquad Q_n=q_n-n,\qquad
 B_n^*=\floor{\varepsilon_nq_n}\ge1,\qquad
 \tau_n=\frac1n\log_2B_n^*.
 \end{gathered}
\]
We abbreviate the corresponding official quantities from
\cref{prob:finite-prize} by
\(\delta_{n,\mathrm{off,sup}}^*\) and
\(\delta_{n,\mathrm{off,grid}}^*\).  Let
\(H_2(x)=-x\log_2x-(1-x)\log_2(1-x)\).  For
\(r\in(0,1)\) and \(b\ge0\), define
\[
 g^*(r,b)=\sup\left(\{g\in[0,1-r]:H_2(r+g)\ge bg\}\cup\{0\}\right),
 \qquad g_n^*=g^*(\rho_n,\beta_n).
\]
The exact locator-prefix mean for
the dimension-\(k_n+1\) lower route is
\[
 \barN_{1,n}(a)=\binom na\abs{\B_n}^{-(a-k_n-1)}.    \tag{6.1}\label{eq:identity-mean}
\]
For an integer \(a\), put \(g=(a-k_n-1)/n\).  Stirling's formula gives,
uniformly away from rates zero and one,
\[
 \frac1n\log_2\barN_{1,n}(a)
 =H_2(\rho_n^++g)-\beta_ng+O\!\left(\frac{\log n}{n}\right). \tag{6.2}\label{eq:identity-exponent}
\]
The pole conversion now imposes an additional exact branch condition.  If
\(Q_n=0\), there is no off-domain pole.  If
\(0<Q_n\le k_nB_n^*\), \cref{lem:exact-pole-conversion} says that the
certified simple-pole lower bound never exceeds \(B_n^*\), regardless of the
list size; this route then supplies no unsafe certificate.  Suppose henceforth that
\(Q_n>k_nB_n^*\), and put
\[
 T_{{\rm pole},n}=T_{\rm pole}(B_n^*;Q_n,k_n),\qquad
 \tau_n^{\rm pole}=\frac1n\log_2T_{{\rm pole},n},\qquad
 A_n=\frac{Q_n-k_n}{Q_n-k_nB_n^*}\ge1.
\]
Since \(B_n^*A_n<T_{{\rm pole},n}\le B_n^*A_n+1\),
\[
 0<\tau_n^{\rm pole}-\tau_n-\frac1n\log_2A_n
 \le\frac1n\log_2\!\left(1+\frac1{B_n^*A_n}\right)
 \le\frac1n.                                      \tag{6.3}\label{eq:pole-exponent-comparison}
\]
Thus \(\tau_n=o(1)\) alone does not control the pole-adjusted level.  The
conditions \(Q_n>k_nB_n^*\), \(\tau_n=o(1)\), and
\(n^{-1}\log_2A_n=o(1)\) imply \(\tau_n^{\rm pole}=o(1)\).  The stronger
collision-nonbinding condition \(k_nB_n^*/Q_n=o(1)\) gives
\(\tau_n^{\rm pole}=\tau_n+O(1/n)\); the error need not be \(o(1/n)\) when
\(B_n^*\) remains bounded.

The exact test for this certified pole lower bound to exceed \(B_n^*\) is
\(\barN_{1,n}(a)>T_{{\rm pole},n}-1\).  Because
\(T_{{\rm pole},n}\ge2\),
\[
 0\le\frac1n\log_2\frac{T_{{\rm pole},n}}
 {T_{{\rm pole},n}-1}\le\frac1n.                 \tag{6.4}\label{eq:pole-ceiling-error}
\]
It is therefore convenient to use the pole-adjusted analytic crossing
\[
 g_{{\rm pole},n}=\sup\left(
 \left\{g\in[0,1-\rho_n^+]:
 H_2(\rho_n^++g)-\beta_ng\ge\tau_n^{\rm pole}\right\}\cup\{0\}
 \right).                                          \tag{6.5}\label{eq:pole-crossing}
\]
For a real \(\xi\), let \(g_{{\rm pole},n}(\xi)\) denote the same right
crossing with \(\tau_n^{\rm pole}\) replaced by
\(\tau_n^{\rm pole}+\xi\).

\begin{lemma}[Stability of the pole-adjusted crossing]\label{lem:crossing-stability}
Assume \(\rho_n\to\rho\in(0,1)\), \(\beta_n\to\infty\),
\(Q_n>k_nB_n^*\), and \(\tau_n^{\rm pole}=o(1)\).  For every sequence
\(0\le\eta_n=o(1)\), uniformly for \(\abs{\xi}\le\eta_n\),
\[
 g_{{\rm pole},n}(\xi)=g_n^*
 +O_\rho\!\left(\frac{\tau_n^{\rm pole}+\eta_n+n^{-1}}{\beta_n}\right)
 =g_n^*+o(g_n^*).                                  \tag{6.6}\label{eq:crossing-stability}
\]
Consequently a uniform multiplicative certificate loss \(2^{\eta_n n}\),
with \(\eta_n=o(1)\), changes the crossing by \(o(g_n^*)\).
\end{lemma}

\begin{proof}
Uniformly for \(r\) in a compact subinterval of \((0,1)\), the positive
right zero of \(H_2(r+g)-\beta g\) is
\(\Theta_\rho(\beta^{-1})\).  On this scale its derivative in \(g\) is
\(-\beta+O_\rho(1)\).  The mean-value theorem shows that changing the right
side from zero to \(\tau_n^{\rm pole}+\xi\) moves the crossing by
\(O_\rho((\tau_n^{\rm pole}+\abs\xi)/\beta_n)\).  Replacing \(\rho_n\) by
\(\rho_n^+=\rho_n+n^{-1}\) moves it by
\(O_\rho((n\beta_n)^{-1})\).  Since
\(g_n^*=\Theta_\rho(\beta_n^{-1})\), the total error is \(o(g_n^*)\).
\end{proof}

\begin{theorem}[Certificate bracket and pole-adjusted crossing]
\label{thm:target-crossing}
Assume \(Q_n>k_nB_n^*\).  Suppose there are agreements
\(a_{-,n}<a_{+,n}\) in \(\{k_n+1,\ldots,n\}\) for which an exact lower certificate proves
\(B_n(a_{-,n})>B_n^*\) and a closed exact upper certificate proves
\(B_n(a_{+,n})\le B_n^*\).  If
\[
 a_{\pm,n}=k_n+1+g_{{\rm pole},n}n+o(n),           \tag{6.7}\label{eq:coarse-bracket}
\]
then the first safe agreement satisfies
\(a_n^*=k_n+1+g_{{\rm pole},n}n+o(n)\), and in a nonsaturated row the
largest safe grid radius and official real supremum are both
\[
 1-\rho_n-g_{{\rm pole},n}+o(1).                  \tag{6.8}\label{eq:target-frontier}
\]
If the two errors in \eqref{eq:coarse-bracket} are \(o(g_n^*n)\),
\(\tau_n^{\rm pole}=o(1)\), and \(\beta_n=o(n)\), then
\[
 g_{{\rm pole},n}=g_n^*+o(g_n^*),\qquad
 \delta_{n,\mathrm{off,sup}}^*=1-\rho_n-g_n^*+o(g_n^*). \tag{6.9}\label{eq:sharp-frontier}
\]
\end{theorem}

\begin{proof}
Monotonicity places the first safe agreement in
\((a_{-,n},a_{+,n}]\).  By \cref{prop:staircase}, in a nonsaturated row
\[
 \delta_{n,\mathrm{off,grid}}^*=1-\frac{a_n^*}{n},
 \qquad
 \delta_{n,\mathrm{off,sup}}^*=1-\frac{a_n^*-1}{n},
\]
which proves \eqref{eq:target-frontier}.  For the sharper statement,
\(D_n\subseteq\B_n\) and \(\abs D_n=n\) imply
\(\beta_n\ge\log_2n\).  Apply \cref{lem:crossing-stability}; since
\(\beta_n=o(n)\), one has \(g_n^*n\to\infty\), so the endpoint discrepancy
\(1/n\) is \(o(g_n^*)\).
\end{proof}

\begin{corollary}[Explicit pole-adjusted crossing]\label{cor:crossing-expansion}
Fix \(0<\rho_0<1/2\), and suppose
\(\rho_n\in[\rho_0,1-\rho_0]\), \(\beta_n\to\infty\), and
\(\tau_n^{\rm pole}=o(1)\).  Uniformly in this range,
\[
 g_n^*=\frac{H_2(\rho_n)}
 {\beta_n-H_2'(\rho_n)}+O_{\rho_0}(\beta_n^{-3}),  \tag{6.10}\label{eq:gstar-expansion}
\]
and
\[
 g_{{\rm pole},n}=\frac{H_2(\rho_n^+)-\tau_n^{\rm pole}}
 {\beta_n-H_2'(\rho_n^+)}
 +O_{\rho_0}(\beta_n^{-3}).                        \tag{6.11}\label{eq:gpole-expansion}
\]
In the collision-nonbinding pole regime, the coefficient-and-folding
field sets the leading displacement from capacity; the challenge field
enters through the target budget and the pole-conversion condition.
\end{corollary}

\begin{proof}
Taylor expansion gives
\(H_2(\rho+g)=H_2(\rho)+H_2'(\rho)g+O_{\rho_0}(g^2)\).
The root equation implies \(g=O(\beta^{-1})\); substituting this estimate
back and dividing by \(\beta-H_2'(\rho)\asymp\beta\) gives
\eqref{eq:gstar-expansion}.  The same calculation at level
\(\tau_n^{\rm pole}\), with \(\rho_n^+\), gives
\eqref{eq:gpole-expansion}.  A sharper expansion may retain the quadratic
entropy term.
\end{proof}

\begin{corollary}[Raw-budget form under subexponential pole conversion]
\label{cor:raw-budget-form}
Suppose \(Q_n>k_nB_n^*\), \(\tau_n=o(1)\),
\(n^{-1}\log_2A_n=o(1)\), and \(\beta_n=o(n)\).  A sharp certificate bracket
as in \cref{thm:target-crossing} then gives
\[
 \delta_{n,\mathrm{off,sup}}^*
 =1-\rho_n-g^*(\rho_n,\log_2\abs{\B_n})+o(g_n^*). \tag{6.12}\label{eq:raw-budget-frontier}
\]
In particular, \(k_nB_n^*/Q_n=o(1)\) is a sufficient collision-nonbinding
hypothesis for the pole-conversion part of this statement.
\end{corollary}

\begin{proof}
Equation~\eqref{eq:pole-exponent-comparison} gives
\(\tau_n^{\rm pole}=o(1)\); now apply
\cref{lem:crossing-stability,thm:target-crossing}.
\end{proof}

The pole condition is a genuine restriction, not a technical restatement of
subexponential budget.  For example, if \(\varepsilon_n=\varepsilon>0\),
\(k_n\sim\rho n\), \(Q_n/q_n\) is bounded away from zero, and
\(B_n^*\sim\varepsilon q_n\), then
\(k_nB_n^*/Q_n=\Theta(n)\); hence pole admissibility eventually fails.
Such a sequence requires a target shrinking on the scale \(1/k_n\) or a
different unsafe mechanism before an identity-pole bracket can be proposed.

Because \(D_n\subseteq\B_n\) has \(n\) distinct points,
\(\beta_n\ge\log_2n\), and
\cref{cor:crossing-expansion} gives \(g_n^*\asymp1/\beta_n\).  When
\(\beta_n=o(n)\), the radius correction tends to zero while
\(g_n^*n\asymp n/\beta_n\) diverges.  A loss \(2^{\eta_n n}\), with
\(\eta_n=o(1)\), shifts the crossing by
\(O(\eta_n/\beta_n)=o(g_n^*)\).  Thus a uniform aggregate \(e^{o(n)}\) loss
is sufficient for the sharper identity displacement.  This fact does not
help with the finite margins in \cref{sec:finite}, where an unspecified
polynomial factor may exceed the entire reserve.  The stronger
\(e^{o(g_n^*n)}\) control supplied by a moving-scale argument gives finer
localization but is not necessary for \eqref{eq:raw-budget-frontier}.

The theorem is a reduction, not the missing upper bound.  A full profile
envelope for a received line is the sum of actual first-match payments,
normalized at their realized images and followed through their actual slope
projections.  The identity-candidate condition is entirely lower-side and
supplies no upper bound.  If a family-uniform closed atlas supplies the
matching upper certificate at \eqref{eq:pole-crossing}, the sequence is
identity-dominant in the sense following \cref{conj:identity-frontier}.
Quotient coefficient fields, partial occupancy, planted factors, image
collapse, and ray saturation may move the crossing; then \(g_n^*\) is only
the identity-profile candidate.

\begin{problem}[Uniform asymptotic closure]\label{prob:asymptotic-closure}
For official full-field positive-density multiplicative and Chebyshev
twin-coset sequences, and for circle presentations proved monomially
equivalent to them, construct an explicit family-uniform, non-oracular,
witness-exhaustive declared atlas and certificate
bracket of width \(o(n)\) around the right crossing of the full
image-normalized envelope.  In a pole-admissible identity-candidate branch
with subexponential raw budget and pole-conversion factor, prove a uniform
aggregate \(2^{o(n)}\) loss and a bracket of width \(o(g_n^*n)\), thereby
establishing identity dominance.  The proof must establish the
Sidon payment of \cref{prob:sidon}, the residual support-to-slope incidence,
the balanced-core payment of \cref{prob:balanced-core}, all remaining
algebraic payments, and a reserve showing that their sum stays below the
target.
\end{problem}

If this problem is solved in the stated pole-admissible identity-candidate window,
\cref{cor:raw-budget-form} gives the precise formula
\(\delta_{n,\mathrm{off,sup}}^*=1-\rho_n-
g^*(\rho_n,\log_2\abs{\B_n})+o(g_n^*)\).  If
\(Q_n=0\) or \(0<Q_n\le k_nB_n^*\), this pole lower route gives no such
conclusion.  A zero
integer budget gives an empty safe set, safety at \(k_n+1\) gives official
threshold \(1\), and a binding tangent profile gives the linear-budget edge;
these branches must be tested first.

\section{Conclusion}\label{sec:conclusion}

The paper separates four stable numerical conjectures from the conjectural
profile architecture that might prove them.  The exact compiler identifies
where primitive realized-image bounds, residual slope or codeword payments,
algebraic routing, and first-match add-back must enter; the incidence and
moment results show why none may be omitted, while the non-oracular payment
condition prevents the architecture from encoding the desired answer.
Known nonidentity attacks show why the identity-candidate hypothesis must
test every explicitly declared nonidentity lower profile, while identity
dominance remains the conclusion to be proved.  The exact pole conversion
makes pole admissibility necessary for its stated unsafe route; another
lower mechanism may behave differently.  No theorem here improves the
safe MCA threshold of \cite{ChoShortening}; proving an MCA inequality in
\cref{conj:benchmark} would give an exact adjacent threshold, while
\cref{prob:asymptotic-closure} would yield the pole-adjusted identity formula
and, under subexponential raw budget and pole-conversion factor, its simpler
\(g^*\)-form.  Any
circle transfer still requires a proved monomial equivalence.

\paragraph{Acknowledgments.}
The author thanks Scott Hughes for contributions to the frontier-adjacent
verification materials and the contributors to the \texttt{rs-mca}
repository for experimental audits and discussions that informed the
finite-ledger program.  The release archive preserves the cited file paths
and pinned commit so that the underlying contribution history remains
auditable.

\paragraph{AI-use disclosure.}
OpenAI ChatGPT/Codex was used substantively to draft and restructure prose,
propose and audit proof steps, check calculations, compare source versions,
and edit \LaTeX.  The author supplied the research program, source
manuscripts, repository materials, mathematical choices, and corrections,
and retains sole responsibility for correctness and originality.

\begin{thebibliography}{BCIKS20/23}
\footnotesize

\bibitem[ABF26]{ABF26}
G.~Arnon, D.~Boneh, and G.~Fenzi,
\newblock Open problems in list decoding and correlated agreement,
\newblock Cryptology ePrint Archive, Paper 2026/680, 2026,
\newblock \url{https://eprint.iacr.org/2026/680}.

\bibitem[BCHKS26]{BCHKS25}
E.~Ben-Sasson, D.~Carmon, U.~Hab\"ock, S.~Kopparty, and S.~Saraf,
\newblock On proximity gaps of Reed--Solomon codes,
\newblock in \emph{Proc. 58th ACM Symposium on Theory of Computing
(STOC '26)}, 2026, 1157--1167,
\newblock \url{https://doi.org/10.1145/3798129.3800827};
extended version ECCC TR25-169 and ePrint 2025/2055.

\bibitem[BCIKS20/23]{BCIKS20}
E.~Ben-Sasson, D.~Carmon, Y.~Ishai, S.~Kopparty, and S.~Saraf,
\newblock Proximity gaps for Reed--Solomon codes,
\newblock \emph{J. ACM} 70 (2023), no.~5, 1--57,
\newblock \url{https://doi.org/10.1145/3614423}.

\bibitem[CS25]{CritesStewart25}
E.~Crites and A.~Stewart,
\newblock On Reed--Solomon proximity gaps conjectures,
\newblock Cryptology ePrint Archive, Paper 2025/2046, 2025,
\newblock \url{https://eprint.iacr.org/2025/2046}.

\bibitem[Cho26a]{ChoShortening}
P.~Chojecki,
\newblock Shortening bounds for Reed--Solomon MCA,
\newblock companion manuscript, version 9.2, July 2026,
\newblock public source at
\href{https://github.com/przchojecki/rs-mca/blob/9908454995f3f195cfe748f35a1135211609d066/experimental/RS_MCA_Paving_v9.2.tex}
{repository snapshot \texttt{9908454995f3}}.

\bibitem[Cho26b]{RSMCARepo}
P.~Chojecki,
\newblock \texttt{rs-mca}: proof development and exact certificate ledgers,
\newblock source repository, pinned snapshot
\href{https://github.com/przchojecki/rs-mca/tree/e190193cebced1d3752d068a1c24136bc69a85d9}{\texttt{e190193cebce}}, 2026,
\newblock \url{https://github.com/przchojecki/rs-mca}, accessed July 19, 2026.

\bibitem[GYXK26]{GaoYangXuKan26}
Y.~Gao, H.~Yang, Y.~Xu, and H.~Kan,
\newblock List-decoding counterexamples yield lower bounds on mutual
correlated agreement error,
\newblock arXiv:2607.10572, 2026,
\newblock \url{https://arxiv.org/abs/2607.10572}.

\bibitem[PP26]{ProximityPrize}
{Ethereum Foundation},
\newblock The Proximity Prize: \$1M Reed--Solomon Challenge,
\newblock preliminary version, 2026,
\newblock \url{https://proximityprize.org/}, accessed July 19, 2026.

\bibitem[Hab25]{Habock25}
U.~Hab\"ock,
\newblock A note on mutual correlated agreement for Reed--Solomon codes,
\newblock Cryptology ePrint Archive, Paper 2025/2110, 2025,
\newblock \url{https://eprint.iacr.org/2025/2110}.

\bibitem[HLP24]{HLP24}
U.~Hab\"ock, D.~Levit, and S.~Papini,
\newblock Circle STARKs,
\newblock Cryptology ePrint Archive, Paper 2024/278, 2024,
\newblock \url{https://eprint.iacr.org/2024/278}.

\bibitem[Jo26]{Jo26}
S.~Jo,
\newblock Reed--Solomon mutual correlated agreement beyond the Johnson radius,
\newblock Cryptology ePrint Archive, Paper 2026/1432, 2026,
\newblock \url{https://eprint.iacr.org/2026/1432}.

\bibitem[KKH26]{KKH26}
D.~Krachun, S.~Kazanin, and U.~Hab\"ock,
\newblock Failure of proximity gaps close to capacity,
\newblock Cryptology ePrint Archive, Paper 2026/782, 2026,
\newblock \url{https://eprint.iacr.org/2026/782}.

\end{thebibliography}

\end{document}
